% Authors: Dr. Denys Dutykh (Mathematics Department, Khalifa University of Science and Technology,
% Abu Dhabi, UAE) and Prof. Laurent Vuillon (Univ. Savoie Mont Blanc, CNRS, LAMA, Chambery, France)
\section{Path envelopes and the first seven rows}\label{sec:early-rows}

The Euclidean bijection turns a cost bound into a finite binary tree.
Enumerating that tree at every denominator would obscure the reason that
an entire row can be treated. Instead, we compare its branches with a
small number of explicit paths. Fibonacci growth controls the cost,
and deviations from the Fibonacci path determine the finer norm bands.

\subsection{Sorted paths and a stable order}

On sorted positive pairs define
\begin{equation}\label{eq:path-maps}
 \mathsf A(x,y)=(x,x+y),\qquad
 \mathsf B(x,y)=(y,x+y).
\end{equation}
The first move from $(1,1)$ is symmetric and gives $(1,2)$. Thereafter
these two moves are distinct, and every sorted primitive pair of cost
$c\geqslant1$ is obtained by a unique sequence of $c-1$ moves.
Indeed, from $0<x<y$ its sorted predecessor is the sorted pair formed
from $x,y-x$. If $y>2x$, the last move was $\mathsf A$; if $y<2x$,
it was $\mathsf B$. The equality $y=2x$ is primitive only at $(1,2)$.
This proves the asserted uniqueness directly, including its initial
boundary.

A move $\mathsf A$ after the fixed first move is called a \emph{defect}.
Write $v_c(D)$ for the endpoint whose defect costs form
$D\subseteq\{2,\ldots,c\}$. Thus a defect at cost $j$ is the move
arriving at depth $j$. Products written as words in the moves are read
chronologically from left to right; an expression $\mathsf B^nz$,
with a vector $z$, means the $n$-fold application of $\mathsf B$ to that vector.

Let $F_0=0,F_1=1,F_{n+1}=F_n+F_{n-1}$, and use
$F_{-n}=(-1)^{n+1}F_n$ when a displayed identity reaches a negative
index. Put
$$
 N(x,y)=x^2+y^2,\qquad
 \Lambda(x,y)=x+\varphi y,\qquad
 \varphi=\frac{1+\sqrt5}{2}.
$$

\begin{lemma}[Maximum, replacement, and propagation]\label{lem:path-orders}
At exact cost $c$, the unique maximum norm is
$F_{2c+3}$, attained at $(F_{c+1},F_{c+2})$.
Replacing any defect by $\mathsf B$ strictly increases the final norm.
Finally, simultaneous weak inequalities in $N$ and $\Lambda$ are
preserved by every common $\mathsf B$-tail.
\end{lemma}

\begin{proof}
Coordinate induction in \eqref{eq:path-maps}, starting from $(1,1)$,
gives $x\leqslant F_{c+1}$ and $y\leqslant F_{c+2}$.
The identity $F_{c+1}^2+F_{c+2}^2=F_{2c+3}$ proves the norm bound.
After the first move, $x<y$, and
$$
 \mathsf B(x,y)-\mathsf A(x,y)=(y-x,0).
$$
This difference is nonnegative and nonzero. Each subsequent matrix
is nonnegative and invertible, so it preserves a nonnegative nonzero
difference. Since both endpoint coordinates are positive, the norm
strictly increases. Replacing every defect therefore proves uniqueness
of the Fibonacci maximum as well.
For propagation, direct calculation gives
\begin{equation}\label{eq:perron-propagation}
 \Lambda(\mathsf Bz)=\varphi\Lambda(z),\qquad
 N(\mathsf Bz)=\varphi^{-2}N(z)+\varphi^{-1}\Lambda(z)^2.
\end{equation}
Both coefficients are positive, and $\Lambda$ is positive on the
states under consideration.
\end{proof}

\begin{lemma}[One-defect minimum]\label{lem:one-defect-minimum}
For $c\geqslant2$, the earliest-defect endpoint
$e_c=\mathsf B^{c-2}(1,3)$ minimizes both $N$ and $\Lambda$ among
one-defect paths, and
\begin{equation}\label{eq:one-defect-minimum}
 N(e_c)=5F_{2c-1}.
\end{equation}
\end{lemma}

\begin{proof}
The assertion starts at $c=2$. An inherited one-defect path is covered
by \eqref{eq:perron-propagation}. The only new path creates its defect
at the last move and has endpoint $b_c=(F_c,F_{c+2})$. Fibonacci
addition gives
$$
 b_c-e_c=F_{c-2}(-1,1),\qquad
 N(b_c)-N(e_c)=4F_{c-2}F_{c-1}.
$$
Its $\Lambda$-difference is $(\varphi-1)F_{c-2}$, so induction closes.
The equality for $N(e_c)$ follows by substitution from
$e_c=(F_c+F_{c-2},F_{c+1}+F_{c-1})$, or from the recurrence
$N(\mathsf B^2z)=3N(\mathsf Bz)-N(z)$ and its first two values.
\end{proof}

\subsection{The second envelope}

\begin{lemma}[The second norm envelope]\label{lem:second-envelope}
At cost $c\geqslant5$, every non-Fibonacci endpoint satisfies
\begin{equation}\label{eq:second-envelope}
 N(x,y)\leqslant E_c:=F_c^2+F_{c+2}^2.
\end{equation}
Equality is attained at $(F_c,F_{c+2})$, and also at $(12,19)$ when
$c=6$; there are no other equality cases.
\end{lemma}

\begin{proof}
Let $j$ be the last defect cost and $r=c-j$. Before that defect,
the coordinate bound in Lemma~\ref{lem:path-orders} gives the upper
state $(F_j,F_{j+1})$. Its final endpoint is
$$
 X_{c,j}=F_c+F_rF_j,\qquad
 Y_{c,j}=F_{c+2}-F_rF_{j-1}.
$$
Any earlier defect makes at least one final coordinate strictly smaller.
Expansion and Cassini's identity give the exact deficit
\begin{align}\label{eq:second-envelope-deficit}
 E_c-X_{c,j}^2-Y_{c,j}^2
 =F_r\bigl[
 2(F_jF_{j-1}+(-1)^j)F_{r+1}
 -F_{j-2}F_{j+1}F_r\bigr].
\end{align}
If $r=0$, this is zero. For $r\geqslant1$ and $j\geqslant4$,
replace $F_{r+1}$ by its lower bound $F_r$; the remaining coefficient is
$$
 2F_{j-1}^2-F_{j-2}^2+2(-1)^j>0.
$$
The case $j=2$ is positive directly. When $j=3$, the bracket in
\eqref{eq:second-envelope-deficit} is $2F_{r+1}-3F_r$.
Here $r\geqslant2$; its values at $r=2,3$ are $1,0$, and for
$r\geqslant4$ it is $F_{r-3}>0$. The zero cases give exactly the
two stated endpoints.
\end{proof}

It follows that an occurrence of cost $L\geqslant6$ and label $m$ is
weakly Euclidean-minimal whenever
\begin{equation}\label{eq:one-hole}
 E_{L-1}<m,\qquad m\ne F_{2L+1}.
\end{equation}
At cost $L-1$ only the Fibonacci maximum can lie above $E_{L-1}$.
Every smaller cost has norm at most
$F_{2L-1}<E_{L-1}$. Thus \eqref{eq:one-hole} excludes every
smaller-cost representation.

\Cref{fig:norm-cloud} shows the situation that this criterion exploits.
Normalized by the Fibonacci maximum, the norms of all pairs of one cost
lie between zero and the second envelope, near $0.829$, except for the
isolated value $1$, so that a hole separates the two. The normalized source labels
of a row tend to the constant $C_P$ computed in \cref{sec:row-labels} below.
For $P\leqslant6$ these limits exceed $1$, and $C_7$ falls in the hole; for
$P\geqslant8$ they fall inside the cloud, where a finer separation is needed,
and that is the business of \cref{sec:tails}.

\begin{figure}[tb]
\centering
% Authors: Dr. Denys Dutykh (Mathematics Department, Khalifa University of Science and
% Technology, Abu Dhabi, UAE) and Prof. Laurent Vuillon (Univ. Savoie Mont Blanc, CNRS, LAMA)
% Generated by supplement/code/make_figures.py from exact data; do not edit by hand.
\begin{tikzpicture}[font=\footnotesize, >={Stealth[length=4pt]}]
\fill[sand] (9.038,0) rectangle (10.9,3.45);
\fill[navy] (0,0) rectangle (0.054,3.277);
\fill[navy] (0.054,0) rectangle (0.109,3.3);
\fill[navy] (0.109,0) rectangle (0.164,3.282);
\fill[navy] (0.164,0) rectangle (0.218,3.259);
\fill[navy] (0.218,0) rectangle (0.273,3.233);
\fill[navy] (0.273,0) rectangle (0.327,3.207);
\fill[navy] (0.327,0) rectangle (0.382,3.186);
\fill[navy] (0.382,0) rectangle (0.436,3.157);
\fill[navy] (0.436,0) rectangle (0.49,3.135);
\fill[navy] (0.49,0) rectangle (0.545,3.11);
\fill[navy] (0.545,0) rectangle (0.6,3.089);
\fill[navy] (0.6,0) rectangle (0.654,3.061);
\fill[navy] (0.654,0) rectangle (0.709,3.04);
\fill[navy] (0.709,0) rectangle (0.763,3.019);
\fill[navy] (0.763,0) rectangle (0.818,2.996);
\fill[navy] (0.818,0) rectangle (0.872,2.972);
\fill[navy] (0.872,0) rectangle (0.927,2.947);
\fill[navy] (0.927,0) rectangle (0.981,2.932);
\fill[navy] (0.981,0) rectangle (1.036,2.903);
\fill[navy] (1.036,0) rectangle (1.09,2.889);
\fill[navy] (1.09,0) rectangle (1.145,2.864);
\fill[navy] (1.145,0) rectangle (1.199,2.847);
\fill[navy] (1.199,0) rectangle (1.254,2.819);
\fill[navy] (1.254,0) rectangle (1.308,2.797);
\fill[navy] (1.308,0) rectangle (1.363,2.779);
\fill[navy] (1.363,0) rectangle (1.417,2.756);
\fill[navy] (1.417,0) rectangle (1.472,2.737);
\fill[navy] (1.472,0) rectangle (1.526,2.714);
\fill[navy] (1.526,0) rectangle (1.581,2.695);
\fill[navy] (1.581,0) rectangle (1.635,2.676);
\fill[navy] (1.635,0) rectangle (1.69,2.651);
\fill[navy] (1.69,0) rectangle (1.744,2.642);
\fill[navy] (1.744,0) rectangle (1.799,2.614);
\fill[navy] (1.799,0) rectangle (1.853,2.59);
\fill[navy] (1.853,0) rectangle (1.907,2.568);
\fill[navy] (1.907,0) rectangle (1.962,2.553);
\fill[navy] (1.962,0) rectangle (2.017,2.536);
\fill[navy] (2.017,0) rectangle (2.071,2.515);
\fill[navy] (2.071,0) rectangle (2.126,2.486);
\fill[navy] (2.126,0) rectangle (2.18,2.466);
\fill[navy] (2.18,0) rectangle (2.235,2.451);
\fill[navy] (2.235,0) rectangle (2.289,2.44);
\fill[navy] (2.289,0) rectangle (2.344,2.399);
\fill[navy] (2.344,0) rectangle (2.398,2.381);
\fill[navy] (2.398,0) rectangle (2.453,2.381);
\fill[navy] (2.453,0) rectangle (2.507,2.353);
\fill[navy] (2.507,0) rectangle (2.562,2.323);
\fill[navy] (2.562,0) rectangle (2.616,2.31);
\fill[navy] (2.616,0) rectangle (2.671,2.289);
\fill[navy] (2.671,0) rectangle (2.725,2.274);
\fill[navy] (2.725,0) rectangle (2.78,2.243);
\fill[navy] (2.78,0) rectangle (2.834,2.236);
\fill[navy] (2.834,0) rectangle (2.889,2.203);
\fill[navy] (2.889,0) rectangle (2.943,2.191);
\fill[navy] (2.943,0) rectangle (2.998,2.178);
\fill[navy] (2.998,0) rectangle (3.052,2.148);
\fill[navy] (3.052,0) rectangle (3.107,2.117);
\fill[navy] (3.107,0) rectangle (3.161,2.109);
\fill[navy] (3.161,0) rectangle (3.216,2.089);
\fill[navy] (3.216,0) rectangle (3.27,2.063);
\fill[navy] (3.27,0) rectangle (3.325,2.047);
\fill[navy] (3.325,0) rectangle (3.379,2.022);
\fill[navy] (3.379,0) rectangle (3.433,2.011);
\fill[navy] (3.433,0) rectangle (3.488,1.969);
\fill[navy] (3.488,0) rectangle (3.543,1.988);
\fill[navy] (3.543,0) rectangle (3.597,1.938);
\fill[navy] (3.597,0) rectangle (3.652,1.943);
\fill[navy] (3.652,0) rectangle (3.706,1.903);
\fill[navy] (3.706,0) rectangle (3.76,1.896);
\fill[navy] (3.76,0) rectangle (3.815,1.859);
\fill[navy] (3.815,0) rectangle (3.869,1.827);
\fill[navy] (3.869,0) rectangle (3.924,1.771);
\fill[navy] (3.924,0) rectangle (3.978,1.8);
\fill[navy] (3.978,0) rectangle (4.033,1.8);
\fill[navy] (4.033,0) rectangle (4.088,1.803);
\fill[navy] (4.088,0) rectangle (4.142,1.766);
\fill[navy] (4.142,0) rectangle (4.197,1.703);
\fill[navy] (4.197,0) rectangle (4.251,1.683);
\fill[navy] (4.251,0) rectangle (4.306,1.688);
\fill[navy] (4.306,0) rectangle (4.36,1.575);
\fill[navy] (4.36,0) rectangle (4.415,1.678);
\fill[navy] (4.415,0) rectangle (4.469,1.531);
\fill[navy] (4.469,0) rectangle (4.524,1.6);
\fill[navy] (4.524,0) rectangle (4.578,1.598);
\fill[navy] (4.578,0) rectangle (4.633,1.523);
\fill[navy] (4.633,0) rectangle (4.687,1.531);
\fill[navy] (4.687,0) rectangle (4.742,1.606);
\fill[navy] (4.742,0) rectangle (4.796,1.464);
\fill[navy] (4.796,0) rectangle (4.851,1.523);
\fill[navy] (4.851,0) rectangle (4.905,1.482);
\fill[navy] (4.905,0) rectangle (4.96,1.421);
\fill[navy] (4.96,0) rectangle (5.014,1.485);
\fill[navy] (5.014,0) rectangle (5.069,1.28);
\fill[navy] (5.069,0) rectangle (5.123,1.491);
\fill[navy] (5.123,0) rectangle (5.178,1.358);
\fill[navy] (5.178,0) rectangle (5.232,1.055);
\fill[navy] (5.232,0) rectangle (5.287,1.219);
\fill[navy] (5.287,0) rectangle (5.341,0.964);
\fill[navy] (5.341,0) rectangle (5.396,0.948);
\fill[navy] (5.396,0) rectangle (5.45,1.244);
\fill[navy] (5.45,0) rectangle (5.505,1.268);
\fill[navy] (5.505,0) rectangle (5.559,1.134);
\fill[navy] (5.559,0) rectangle (5.614,1.225);
\fill[navy] (5.614,0) rectangle (5.668,1.331);
\fill[navy] (5.668,0) rectangle (5.723,1.183);
\fill[navy] (5.723,0) rectangle (5.777,1.076);
\fill[navy] (5.777,0) rectangle (5.832,1.143);
\fill[navy] (5.832,0) rectangle (5.886,1.205);
\fill[navy] (5.886,0) rectangle (5.941,0.978);
\fill[navy] (5.941,0) rectangle (5.995,1.134);
\fill[navy] (5.995,0) rectangle (6.05,1.311);
\fill[navy] (6.05,0) rectangle (6.104,0.812);
\fill[navy] (6.104,0) rectangle (6.158,0.786);
\fill[navy] (6.158,0) rectangle (6.213,1.031);
\fill[navy] (6.213,0) rectangle (6.268,0.729);
\fill[navy] (6.268,0) rectangle (6.322,0.219);
\fill[navy] (6.322,0) rectangle (6.377,0.992);
\fill[navy] (6.377,0) rectangle (6.431,0.509);
\fill[navy] (6.431,0) rectangle (6.486,0.219);
\fill[navy] (6.486,0) rectangle (6.54,0.219);
\fill[navy] (6.54,0) rectangle (6.595,0.348);
\fill[navy] (6.649,0) rectangle (6.704,0.219);
\fill[navy] (6.758,0) rectangle (6.812,0.759);
\fill[navy] (6.812,0) rectangle (6.867,0.439);
\fill[navy] (6.867,0) rectangle (6.921,0.759);
\fill[navy] (6.921,0) rectangle (6.976,0.992);
\fill[navy] (6.976,0) rectangle (7.031,1.125);
\fill[navy] (7.031,0) rectangle (7.085,0.835);
\fill[navy] (7.085,0) rectangle (7.14,0.348);
\fill[navy] (7.14,0) rectangle (7.194,0.219);
\fill[navy] (7.194,0) rectangle (7.249,0.964);
\fill[navy] (7.249,0) rectangle (7.303,0.509);
\fill[navy] (7.412,0) rectangle (7.467,0.219);
\fill[navy] (7.467,0) rectangle (7.521,0.812);
\fill[navy] (7.521,0) rectangle (7.575,0.348);
\fill[navy] (7.63,0) rectangle (7.684,0.219);
\fill[navy] (7.684,0) rectangle (7.739,0.219);
\fill[navy] (7.902,0) rectangle (7.957,0.219);
\fill[navy] (8.556,0) rectangle (8.611,0.219);
\fill[navy] (8.666,0) rectangle (8.72,0.567);
\fill[navy] (8.72,0) rectangle (8.775,0.759);
\fill[navy] (8.992,0) rectangle (9.047,0.348);
\draw[wine, line width=1.4pt] (10.9,0) -- (10.9,0.219);
\fill[wine] (10.9,0.219) circle (1.6pt);
\node[wine, align=center, font=\scriptsize] (fib) at (9.969,1.386) {Fibonacci\\ maximum};
\draw[slate, line width=0.4pt] (fib.south) -- (10.86,0.279);
\node[slate, align=center, font=\scriptsize] at (9.969,2.805) {hole};
\draw[slate, densely dashed, line width=0.6pt] (7.951,0) -- (7.951,3.45);
\node[slate, anchor=south east, align=right, font=\scriptsize] at (7.911,3.47) {one-defect\\ minimum};
\node[slate, anchor=south west, align=left, font=\scriptsize] at (9.078,3.47) {second\\ envelope};
\draw[ink, line width=0.5pt] (0,0) -- (10.9,0);
\draw[ink, line width=0.5pt] (0,0) -- (0,-0.08);
\node[ink, below] at (0,-0.08) {$0.0$};
\draw[ink, line width=0.5pt] (1.09,0) -- (1.09,-0.08);
\node[ink, below] at (1.09,-0.08) {$0.1$};
\draw[ink, line width=0.5pt] (2.18,0) -- (2.18,-0.08);
\node[ink, below] at (2.18,-0.08) {$0.2$};
\draw[ink, line width=0.5pt] (3.27,0) -- (3.27,-0.08);
\node[ink, below] at (3.27,-0.08) {$0.3$};
\draw[ink, line width=0.5pt] (4.36,0) -- (4.36,-0.08);
\node[ink, below] at (4.36,-0.08) {$0.4$};
\draw[ink, line width=0.5pt] (5.45,0) -- (5.45,-0.08);
\node[ink, below] at (5.45,-0.08) {$0.5$};
\draw[ink, line width=0.5pt] (6.54,0) -- (6.54,-0.08);
\node[ink, below] at (6.54,-0.08) {$0.6$};
\draw[ink, line width=0.5pt] (7.63,0) -- (7.63,-0.08);
\node[ink, below] at (7.63,-0.08) {$0.7$};
\draw[ink, line width=0.5pt] (8.72,0) -- (8.72,-0.08);
\node[ink, below] at (8.72,-0.08) {$0.8$};
\draw[ink, line width=0.5pt] (9.81,0) -- (9.81,-0.08);
\node[ink, below] at (9.81,-0.08) {$0.9$};
\draw[ink, line width=0.5pt] (10.9,0) -- (10.9,-0.08);
\node[ink, below] at (10.9,-0.08) {$1.0$};
\node[ink] at (5.45,-0.75) {normalized norm $N/F_{43}$ of the pairs of cost $20$};
\draw[mist, line width=0.5pt] (0,-1.35) -- (10.9,-1.35);
\draw[teal, line width=1.1pt] (9.275,-1.47) -- (9.275,-1.23);
\node[teal, below, font=\scriptsize] at (9.275,-1.47) {$C_{7}$};
\draw[wine, line width=1.1pt] (7.69,-1.47) -- (7.69,-1.23);
\node[wine, below, font=\scriptsize] at (7.69,-1.47) {$C_{8}$};
\draw[wine, line width=1.1pt] (6.377,-1.47) -- (6.377,-1.23);
\node[wine, below, font=\scriptsize] at (6.377,-1.47) {$C_{9}$};
\draw[wine, line width=1.1pt] (5.287,-1.47) -- (5.287,-1.23);
\node[wine, below, font=\scriptsize] at (5.287,-1.47) {$C_{10}$};
\draw[wine, line width=1.1pt] (4.384,-1.47) -- (4.384,-1.23);
\node[wine, below, font=\scriptsize] at (4.384,-1.47) {$C_{11}$};
\draw[wine, line width=1.1pt] (3.635,-1.47) -- (3.635,-1.23);
\node[wine, below, font=\scriptsize] at (3.635,-1.47) {$C_{12}$};
\draw[wine, line width=1.1pt] (3.014,-1.47) -- (3.014,-1.23);
\node[wine, below, font=\scriptsize] at (3.014,-1.47) {$C_{13}$};
\draw[teal, ->, line width=0.7pt] (10.95,-1.35) -- (11.65,-1.35);
\node[teal, right, font=\scriptsize, align=left] at (11.65,-1.35) {$C_2,\ldots,C_6$\\ above $1$};
\node[ink, left, font=\scriptsize] at (-0.05,-1.35) {sources};
\node[ink, rotate=90, font=\scriptsize] at (-0.45,1.65) {pairs (log scale)};
\end{tikzpicture}

\caption{The normalized norms $N/F_{43}$ of all $2^{19}$ sorted primitive pairs
of cost $20$, counted in bins of width $1/200$ on a logarithmic scale (navy).
The Fibonacci maximum is the isolated value $1$; every other pair lies at or
below the second envelope of \cref{lem:second-envelope}, and every
one-defect pair at or above the minimum of \cref{lem:one-defect-minimum}
(dashed). The shaded hole contains no normalized norm. The ticks below mark
the limits $C_P$ of the normalized source labels of the rows: $C_7$ lies in
the hole, $C_8,\ldots,C_{13}$ inside the cloud, and $C_2,\ldots,C_6$ above $1$.}
\label{fig:norm-cloud}
\end{figure}

\subsection{The labels of a row}\label{sec:row-labels}

The source labels of the later rows must be placed in narrow normalized
intervals, uniformly in the denominator. This subsection computes the
constant $C_P$ that they approach and bounds the approach on both sides,
with explicit constants. The idea is to diagonalize the Cohn matrix
$\mu(a)$. In its eigenbasis the second Cohn matrix $\mu(b)$ becomes a
matrix with positive entries, so that the trace formula
\eqref{eq:cohn-source} turns the label into a sum of positive terms; the
term that follows the dominant eigenvector at every letter gives $C_P$,
and the others are corrections whose size is governed by the runs of the
letter $a$.

Put
$$
 \lambda=\frac3{\sqrt5\,\varphi}=\frac{15-3\sqrt5}{10},\qquad
 C_P=\varphi^2\lambda^{P-1},\qquad
 \widehat B=\frac1{\sqrt5}\begin{pmatrix}3\varphi^3&2\\2&3\varphi^{-3}\end{pmatrix},
$$
and, for a reduced fraction $P/Q$, write its lower Christoffel word as
$$
 \gamma_{P/Q}=a^{u_1}b\,a^{u_2}b\cdots a^{u_P}b,\qquad u_k\geqslant0 .
$$
The normalized label of the occurrence is
\begin{equation}\label{eq:row-ratio}
 \rho_{P,Q}=\frac{m_{P/Q}}{F_{2(P+Q-2)+1}} ,
\end{equation}
the label divided by the Fibonacci maximum one cost below the occurrence
cost.

\begin{proposition}[The labels of a row]\label{prop:row-labels}
Let $0<P<Q$ be coprime and $n=\lfloor(Q-P)/P\rfloor$.
\begin{enumerate}[label=\textup{(\roman*)},ref=\textup{(\roman*)},leftmargin=*]
\item\label{it:row-runs} Every run satisfies $n\leqslant u_k\leqslant n+1$, and $u_1+\cdots+u_P=Q-P$.
\item\label{it:row-trace} The label is the trace
 $$
 m_{P/Q}=\frac13\tr\prod_{k=1}^{P}
 \begin{pmatrix}\varphi^{2u_k}&0\\0&\varphi^{-2u_k}\end{pmatrix}\widehat B ,
 $$
 the factors being multiplied in the order $k=1,\dots,P$.
\item\label{it:row-bounds} If $P\geqslant2$, then
 \begin{equation}\label{eq:row-source-limit}
 C_P<\rho_{P,Q}<C_P\bigl(1+\varphi^{-4n-6}\bigr)^P ,
 \end{equation}
 while $\rho_{1,Q}<C_1$ for the first row. In particular
 $\rho_{P,Q}$ tends to $C_P$ when $Q$ grows with $P$ fixed.
\end{enumerate}
\end{proposition}

\begin{proof}
\emph{The runs.} The $k$th letter $b$ of $\gamma_{P/Q}$ stands at position
$\lceil kQ/P\rceil$, as recalled after \cref{prop:occurrence-cost}, so that
$u_1=\lceil Q/P\rceil-1$ and
$u_k=\lceil kQ/P\rceil-\lceil(k-1)Q/P\rceil-1$ for $k\geqslant2$. For every real
$x$ and $y$ the difference $\lceil x+y\rceil-\lceil x\rceil$ is $\lfloor y\rfloor$
or $\lceil y\rceil$; with $y=Q/P$, every $u_k$ is therefore
$\lfloor Q/P\rfloor-1=n$ or $\lceil Q/P\rceil-1\leqslant n+1$. The word has $Q$
letters, of which $P$ are $b$, so the runs add up to $Q-P$.

\emph{The eigenbasis.} Let $V$ be the matrix with columns
$(\varphi,1)^{\mathsf T}$ and $(1,-\varphi)^{\mathsf T}$. Its columns are
orthogonal, each of squared length $\varphi^2+1=\sqrt5\,\varphi$, so that
$O=V/(5^{1/4}\varphi^{1/2})$ is an orthogonal matrix. Since
$\varphi^2=\varphi+1$,
$$
 \mu(a)\binom\varphi1=\binom{2\varphi+1}{\varphi+1}=\varphi^2\binom\varphi1,
 \qquad
 \mu(a)\binom1{-\varphi}=\binom{2-\varphi}{1-\varphi}=\varphi^{-2}\binom1{-\varphi},
$$
because $\varphi^3=2\varphi+1$ and $\varphi^{-2}=2-\varphi$. Thus
$O^{\mathsf T}\mu(a)O=\operatorname{diag}(\varphi^2,\varphi^{-2})$. For the second
matrix,
$$
 V^{\mathsf T}\mu(b)V=
 \begin{pmatrix}5\varphi^2+4\varphi+1&2\varphi\\2\varphi&6-3\varphi\end{pmatrix}
 =\begin{pmatrix}3\varphi^4&2\varphi\\2\varphi&3\varphi^{-2}\end{pmatrix},
$$
using $5\varphi^2+4\varphi+1=9\varphi+6=3\varphi^4$ and $6-3\varphi=3\varphi^{-2}$;
dividing by $\sqrt5\,\varphi$ gives $O^{\mathsf T}\mu(b)O=\widehat B$. All four
entries of $\widehat B$ are positive, $\tr\widehat B=6$ and $\det\widehat B=1$.
Since $\mu(\gamma_{P/Q})=\prod_k\mu(a)^{u_k}\mu(b)$ and
$O^{\mathsf T}\mu(a)^{u}O=\operatorname{diag}(\varphi^{2u},\varphi^{-2u})$, conjugating
by $O$ and taking traces turns \eqref{eq:cohn-source} into \ref{it:row-trace}.

\emph{Normalization.} Write $\beta=\widehat B_{12}/\widehat B_{11}=2/(3\varphi^3)$,
$\delta=\widehat B_{22}/\widehat B_{11}=\varphi^{-6}$ and $t_k=\varphi^{-4u_k}$. Taking
the factor $\varphi^{2u_k}\widehat B_{11}$ out of the $k$th matrix of \ref{it:row-trace} leaves
$$
 M_k=\begin{pmatrix}1&\beta\\t_k\beta&t_k\delta\end{pmatrix},\qquad
 m_{P/Q}=\frac13\,\varphi^{2(Q-P)}\Bigl(\frac{3\varphi^3}{\sqrt5}\Bigr)^P\Theta,
 \qquad \Theta=\tr\prod_{k=1}^PM_k ,
$$
by \ref{it:row-runs}. At the odd index $2(P+Q)-3$, Binet's formula reads
$F_{2(P+Q)-3}=\bigl(\varphi^{2(P+Q)-3}+\varphi^{3-2(P+Q)}\bigr)/\sqrt5$. Dividing,
and collecting the powers of $3$, $\sqrt5$ and $\varphi$,
$$
 \rho_{P,Q}
 =\frac{3^{P-1}\varphi^{3-P}}{5^{(P-1)/2}}\cdot\frac{\Theta}{1+\eta}
 =C_P\,\frac{\Theta}{1+\eta},\qquad \eta=\varphi^{6-4(P+Q)} .
$$

\emph{The lower bound.} Expand the trace of the product of the
nonnegative matrices $M_k$ as a sum over closed walks
$i_1,\dots,i_P,i_{P+1}=i_1$ in $\{1,2\}$, the walk contributing
$\prod_k(M_k)_{i_ki_{k+1}}$. Every contribution is nonnegative. The walk that
stays at $1$ contributes $1$. When $P\geqslant2$, the walk that visits $2$
only at the step $j=1$ contributes $(M_P)_{12}(M_1)_{21}=\beta^2t_1$. Hence
$\Theta\geqslant1+\beta^2t_1\geqslant1+\beta^2\varphi^{-4(Q-P)}$, since $u_1\leqslant Q-P$.
The ratio of $\beta^2\varphi^{-4(Q-P)}$ to $\eta$ is
$$
 \frac{4}{9\varphi^6}\,\varphi^{-4Q+4P}\,\varphi^{4P+4Q-6}
 =\frac49\,\varphi^{8P-12}\geqslant\frac49\,\varphi^4>1 ,
$$
because $P\geqslant2$ and $\varphi^4=(7+3\sqrt5)/2>9/4$. Therefore $\Theta>1+\eta$
and $\rho_{P,Q}>C_P$.

\emph{The upper bound.} By \ref{it:row-runs}, $0\leqslant t_k\leqslant\tau:=\varphi^{-4n}$, so that
each $M_k$ is bounded entrywise by
$M=\left(\begin{smallmatrix}1&\beta\\\tau\beta&\tau\delta\end{smallmatrix}\right)$,
and since products of nonnegative matrices are monotone entrywise,
$\Theta\leqslant\tr M^P$. The matrix $M$ has trace $1+\tau\delta>0$ and determinant
$\tau(\delta-\beta^2)=5\tau/(9\varphi^6)>0$; its discriminant
$(1+\tau\delta)^2-\tfrac{20}{9}\tau\delta=1-\tfrac29\tau\delta+(\tau\delta)^2$ is positive because
$\tau\delta\leqslant\varphi^{-6}<1$.
Its eigenvalues $\mu_1,\mu_2$ are therefore positive, and
$$
 \tr M^P=\mu_1^P+\mu_2^P\leqslant(\mu_1+\mu_2)^P=(1+\tau\delta)^P
 =\bigl(1+\varphi^{-4n-6}\bigr)^P ,
$$
the binomial expansion of the right-hand side containing the left-hand
side among nonnegative terms. As $1+\eta>1$, the upper bound in
\eqref{eq:row-source-limit} follows.

\emph{The first row.} When $P=1$ the product has the single factor $M_1$ and
$u_1=Q-1$, so that $\Theta=1+t_1\delta=1+\varphi^{-4Q-2}$, which is smaller than
$1+\eta=1+\varphi^{2-4Q}$. Hence $\rho_{1,Q}<C_1$, and $\rho_{1,Q}=C_1(1+\varphi^{-4Q-2})/(1+\varphi^{2-4Q})$ tends to
$C_1$. Finally, when $P\geqslant2$, the integer $n\geqslant(Q-2P+1)/P$ grows without bound
with $Q$, and both sides of \eqref{eq:row-source-limit} tend to $C_P$.
\end{proof}

The proposition shows where the constant comes from. In the leading
term of \ref{it:row-trace}, each letter $a$ contributes the factor $\varphi^2$ and each
letter $b$ the factor $\widehat B_{11}=3\varphi^3/\sqrt5=\lambda\varphi^4$; the
Fibonacci normalization grows by $\varphi^2$ for each unit of cost, a letter $a$
adding one unit and a letter $b$ two. Every letter $b$ therefore costs
exactly the factor $\lambda$ against the Fibonacci growth, and
$C_P=(\varphi^2/\lambda)\lambda^P$. The same number appears on the side of the
paths. Normalized by the maximum at its own cost, the second envelope of
\cref{lem:second-envelope} tends to $\lambda$, since
$E_c/F_{2c+3}=(F_c^2+F_{c+2}^2)/(F_{c+1}^2+F_{c+2}^2)$ and $1+\varphi^4=3\varphi^2$;
and the second row lies on that envelope, because for odd $Q=2n+3$ the
word $\gamma_{2/Q}$ has central factor $a^nba^n$, whose Gram factor
$P_2\,\mu(a)^n$ sends $e=(F_3,F_2)^{\mathsf T}$ to $(F_{Q+2},F_Q)^{\mathsf T}$, so
that $m_{2/Q}=E_Q$. The value $\lambda$ is also the limit of the one-defect
ratio of \cref{fig:parity-order}(a).

Lee, Li, Rabideau and Schiffler prove, in their coordinates, that for fixed
$P\geqslant2$ the ratio \eqref{eq:row-ratio} does not increase with $Q$
\cite[Lemma~6.4]{LeeEtAl2023}, and they compute its limit in the proof of
their Lemma~6.6 \cite[equation~(6.4)]{LeeEtAl2023}.\footnote{Two printed slips
of that proof should be corrected. The equality labels the first-row value
$m_{q,1}$ as $F_{2q-1}$ on the line that carries the Binet expression of
$F_{2q+1}$, which is the correct index. And the monotonicity is asserted there
for every positive parameter, whereas at the parameter of the first row the
ratio is $F_{2Q+1}/F_{2Q-1}$, which increases strictly to $\varphi^2$ and stays
below its limit, as \cref{prop:row-labels}\ref{it:row-bounds} also shows.}
Neither result is used below: \cref{prop:row-labels} gives the lower bound
for every denominator directly, and its upper bound, completed by finitely
many exact comparisons, replaces the monotonicity in \cref{sec:assembly}.

\subsection{The first seven rows}

With the second envelope and the labels of a row at hand, the first seven
rows follow from norms alone: their source labels sit above every norm of
smaller cost except possibly the Fibonacci maximum, and that maximum is
avoided.

\begin{proposition}\label{prop:rows-seven}
Every occurrence with $1\leqslant P\leqslant7$ is weakly
Euclidean-minimal.
\end{proposition}

\begin{proof}
The first row has $m_{1/Q}=F_{2Q+1}$ and cost $L=Q-1$.
Every smaller cost has norm at most $F_{2Q-1}$, proving this case.
The row formula follows, for example, by iterating the first Cohn
matrix in \eqref{eq:cohn-source}.

For $L\geqslant6$ put $u=F_{L-1}/F_L$. The interval
$3/5\leqslant u\leqslant5/8$ is preserved by $u\mapsto1/(1+u)$.
The inequality $E_{L-1}<(5/6)F_{2L+1}$ reduces to
$7u^2+2u-4<0$, whose maximum on this interval is $-1/64$.
The constants $C_P$ decrease with $P$, and
$$
 C_7=\frac{729}{125\varphi^4}>\frac56,\qquad
 C_6=\frac{243}{25\sqrt5\varphi^3}>1.
$$
The two radical signs reduce respectively to
$4373>1875\sqrt5$ and $59>25\sqrt5$, with squared margins
$1545004$ and $356$. Hence
\eqref{eq:row-source-limit} proves \eqref{eq:one-hole} for
$2\leqslant P\leqslant6$. For $P=7$, the lower bound of
\eqref{eq:row-source-limit} and $C_7>5/6$ give $m>E_{L-1}$, and it remains to
show that $m\neq F_{2L+1}$, that is $\rho_{7,Q}\neq1$; we show $\rho_{7,Q}<1$. When
$Q\geqslant14$, every run of $\gamma_{7/Q}$ has length at least
$\lfloor(Q-7)/7\rfloor\geqslant1$, and \cref{prop:row-labels} gives
$\rho_{7,Q}<C_7(1+\varphi^{-10})^7$. Since $\varphi^{10}=(123+55\sqrt5)/2>120$ and
$\varphi^4=(7+3\sqrt5)/2>6.85$, one has
$(1+\varphi^{-10})^7<e^{7/120}<e^{0.06}<1.07$ and
$C_7=729/(125\varphi^4)<729/856<0.86$, so that $\rho_{7,Q}<0.86\cdot1.07<0.93$.
For $8\leqslant Q\leqslant13$ the six labels
$195025$, $499393$, $1278818$, $3276509$, $8399329$ and $21531778$ lie below
$F_{27}=196418$, $F_{29}=514229$, $F_{31}=1346269$, $F_{33}=3524578$,
$F_{35}=9227465$ and $F_{37}=24157817$ respectively.
Thus this row also avoids the Fibonacci hole.

The remaining $L\leqslant5$ slopes are
$1/2,1/3,1/4,1/5,1/6,2/3,2/5,3/4$.
Their $(L,m)$ values are respectively
$$
 (1,5),(2,13),(3,34),(4,89),(5,233),(3,29),(5,194),(5,169).
$$
Each satisfies $m>F_{2L+1}$, the maximum at costs at most $L-1$.
\end{proof}

The later rows require a finer separation than the second envelope.
The next section classifies fixed-defect paths by their gap parities, and
Section~\ref{sec:assembly} combines those unbounded covers with the
complete finite Gaussian fibers of Section~\ref{sec:gaussian}.
